% =====================================================================
\section{ETSI GS QKD 018}
\label{sec:qkd018}
% =====================================================================

GS~QKD~018 (V1.1.1, April 2022)~\cite{qkd018} defines the
\emph{orchestration interface} between an SDN orchestrator and the SDN
controller of a QKD network. The controller acts as \emph{server}, the
orchestrator as \emph{client}; the information model is expressed in YANG
and is vendor-independent. In SDN terms it is the \emph{northbound
interface} of the QKD controller, i.e.\ the exact complement of
GS~QKD~015. The overall hierarchy becomes:
\begin{center}
\small
orchestrator $\xrightarrow{\text{QKD 018}}$ controller
$\xrightarrow{\text{QKD 015}}$ SDN agent / SD-QKD node
$\xrightarrow{\text{QKD 004/014}}$ SAE.
\end{center}

\begin{definition}[SDN orchestration]
\emph{SDN orchestration} is the continuous process by which available
resources are automatically coordinated, according to optimisation
criteria, to establish and release end-to-end services that span several
network domains, each governed by its own SDN
controller~\cite[§4.2]{qkd018}.
\end{definition}

\begin{definition}[Service link]
A \emph{service link} is a logical key association link between two QKD
nodes, realised by one or more physical QKD links (single-hop or
multi-hop). At the orchestration level it corresponds to the virtual link
of GS~QKD~015.
\end{definition}

% Figure: GS QKD 018 reference scenario (adapted from GS QKD 018, Fig. 1)
\begin{figure}[t]
\centering
\begin{tikzpicture}
  \node[box,fill=blue!8,minimum width=3.3cm] (orch) {SDN orchestrator};
  \node[box,fill=green!8,minimum width=3.3cm,below left=12mm and 6mm of orch] (cq)
       {SDN controller\\QKD network};
  \node[box,fill=orange!10,minimum width=3.3cm,below right=12mm and 6mm of orch] (co)
       {SDN controller\\OTN};
  \node[box,fill=green!4,below=10mm of cq,minimum width=4.6cm] (nq)
       {SD-QKD nodes (agent, KME)};
  \node[box,fill=orange!5,below=10mm of co,minimum width=4.6cm] (no)
       {OTN nodes (SAE)};
  \draw[<->,very thick,red!70!black] (orch) -- node[lbl]{GS QKD 018} (cq);
  \draw[<->,thick] (orch) -- node[lbl]{OTN NBI (out of scope)} (co);
  \draw[<->,thick] (cq) -- node[lbl]{GS QKD 015} (nq);
  \draw[<->,thick] (co) -- (no);
  \draw[<->,dashed] (nq) -- node[lbl]{QKD 004 / 014} (no);
\end{tikzpicture}
\caption{GS~QKD~018 reference scenario (adapted
from~\cite[Fig.~1]{qkd018}). The standardised interface is in red; the
dashed line is key delivery from KMEs to SAEs, which happens inside the
same security perimeter.}
\label{fig:scenario}
\end{figure}


% ---------------------------------------------------------------------
\subsection{GS QKD 018 use cases}
\label{sec:usecase018}

The use case behind the standard (§4.1) takes up almost literally the
backbone problem of GS~QKD~002. An operator introducing QKD can either let
quantum and classical channels coexist in the same fibre or keep them
separate. Coexistence is hard for two reasons: the optical powers of
classical channels are many orders of magnitude higher and disturb the
quantum channel, and the reachable distances of the two kinds of channel
are very different. In a fibre-rich environment the operator may therefore
choose a QKD network separate from the optical transport network (OTN),
gaining performance, isolation and separation of responsibilities.
Following the SDN paradigm, the two networks are governed by different
controllers.

This creates a coordination problem: the SAEs that need to encrypt traffic
sit in OTN nodes, keys are generated in QKD nodes, and each controller knows
only its own network. The orchestrator solves three issues:

\begin{enumerate}
  \item \textbf{Address matching}: it pairs the address of the OTN node that
        must receive keys with that of the QKD node that supplies them; the
        two are inside the same security perimeter (Fig.~\ref{fig:scenario});
  \item \textbf{end-to-end provisioning}: it translates a customer request
        into the configuration of the individual domains, chooses the domains
        to traverse and coordinates the paths;
  \item \textbf{management}: it collects the topology, key-resource
        inventory and FCAPS information of each domain.
\end{enumerate}

Generalising, the same scheme makes it possible to orchestrate
\emph{several QKD domains from different vendors}, each behind its own
controller, as well as QKD and classical domains together. The standard
also points out a practical advantage: an operator that already has an
orchestrator for its optical networks can integrate the QKD network without
introducing separate management tools, which is exactly the integrated
management requirement of UC4.

The functions required of the orchestrator (§4.2) are:
\begin{itemize}
  \item translating customer service requests into domain configuration and
        SAE allocation;
  \item requesting establishment and release of services, including
        inter-domain connections;
  \item computing multi-domain paths, with inter-domain connections and
        endpoints;
  \item changing paths under operator constraints;
  \item discovering the nodes and links of each controller and keeping an
        abstract view of them;
  \item monitoring the key resources available on nodes and links;
  \item monitoring FCAPS management and receiving its notifications.
\end{itemize}

% ---------------------------------------------------------------------
\subsection{The information model}
\label{sec:infomodel}

Clause 5 translates these functions into five information blocks.

\paragraph{Topology discovery (§5.1).}
The controller exposes SD-QKD nodes, interfaces and physical links, and
later service links. The orchestrator does not need the full physical
composition, but an \emph{abstract} view of each domain and of the
connectivity between domains. Discovery can be \emph{proactive} (the
orchestrator queries the controller before every new path computation) or
\emph{reactive} (the controller notifies every change, e.g.\ a node added
or removed).

\paragraph{Network status and inventory (§5.2).}
For each node and physical link the controller exposes status, physical
performance and above all key resources: the generation rate and the rate
available after internal consumption. Service links add bandwidth,
available key and the restoration policy in case of failure.

\paragraph{End-to-end service status (§5.3).}
The orchestrator can query the active external applications, the local and
remote nodes and the service links that serve them, with their status and
key rate.

\paragraph{Provisioning with path computation (§5.4).}
The orchestrator asks the controller for paths that satisfy given
constraints; the controller returns a list of candidates, possibly several
times with different constraints, until the orchestrator picks one to
deploy. In case of failure the orchestrator can request the path to be
replaced.

\paragraph{Notifications (§5.5).}
The controller reports events, alarms and threshold crossings. Performance
thresholds (QBER, SKR, ESKR) are configurable by the orchestrator. Compared
with 015, notifications are added for node status (\texttt{sdqkdn\_new},
\texttt{\_operating}, \texttt{\_down}, \texttt{\_failure},
\texttt{\_out}), for deployed applications
(\texttt{sdqkdn\_application\_deployed}) and for service links
(\texttt{sdqkdn\_svc\_link\_down}, \texttt{\_perf\_update},
\texttt{\_overloaded}).

Equation~\eqref{eq:eskr015} extends naturally to the orchestration level.
Since in a multi-hop path $p$ every link relays the same key, a new service
with demand $r$ is sustainable only if the bottleneck can carry it:
\begin{equation}
  p \text{ admissible} \iff
  \min_{\ell\in p}\, \mathrm{ESKR}_\ell \;\ge\; r .
  \label{eq:admissible}
\end{equation}
This too is our own formalisation, consistent with the \texttt{skr},
\texttt{eskr} and \texttt{expected\_consumption} fields of the YANG models;
the standard does not prescribe any path-computation algorithm.
